\section{Related Work}
\label{sec:related_work}

The architecture of Hexanna IIVX intersects four established research domains: autonomous Text-to-SQL language agents, vectorized columnar analytical engines, out-of-core query evaluation systems, and LLM execution sandboxing.

\subsection{Modern Text-to-SQL and LLM Database Agents}
Translating natural language inquiries into relational database queries has evolved from classical semantic parsing architectures \cite{qin2024survey} to deep neural sequence-to-sequence models and modern Large Language Models \cite{zhou2022table,zan2023large,chen2021codex}. The introduction of cross-domain benchmarks such as Spider 2.0 \cite{zhang2024spider2} and the large-scale, database-grounded BIRD benchmark \cite{li2024bird} highlighted the challenges of schema linking, multi-table joins, and complex relational algebra for neural models.

Recent research has focused on structured prompt engineering and multi-stage agentic workflows. DIN-SQL \cite{pourreza2024dinsql} introduced decomposed in-context learning with self-correction mechanisms, breaking complex SQL synthesis into schema linking, query classification, and syntactic refinement. DAIL-SQL \cite{gao2023dailsql} systematically evaluated candidate selection and prompt representations across LLMs, demonstrating that compact schema representation and few-shot exemplar selection significantly enhance semantic parsing accuracy. MAC-SQL \cite{wang2024macsql} extended this paradigm to multi-agent collaboration, deploying specialized agent roles for decomposition, query generation, and candidate validation. More broadly, iterative reasoning frameworks such as ReAct \cite{yao2022react}, Reflexion \cite{shinn2023reflexion}, Chameleon \cite{lu2023chameleon}, and AutoGen \cite{wu2023autogen} establish that autonomous loop feedback improves generation fidelity across general software and analytical benchmarks \cite{liu2023agentbench}.

However, existing Text-to-SQL research primarily focuses on syntactic and semantic validity over modest-sized, memory-resident relational databases. These frameworks implicitly assume that query execution occurs within managed database servers, failing to address the acute local resource constraints (RAM exhaustion, kernel OOM termination, and context window limits) encountered when autonomous agents execute directly on raw tabular files on commodity developer machines.

\subsection{Vectorized and Columnar Analytical Engines}
Columnar database architectures \cite{idreos2022columnar} achieve superior analytical throughput over row-oriented row stores by storing contiguous values of identical attributes, enabling high compression ratios and efficient CPU cache utilization. Modern vectorized interpretation, systematically benchmarked across cloud and analytical systems \cite{kersten2022vectorized}, processes tabular data in vector chunks (typically 1,024 records), utilizing SIMD vector instructions and avoiding the tuple-at-a-time interpretative overhead of classic database engines.

The open-source ecosystem has witnessed a proliferation of embedded and distributed columnar engines. DuckDB \cite{raasveldt2023duckdb} established an embedded vectorized analytical engine capable of executing complex analytical queries directly against Parquet, Arrow, and CSV formats with zero external dependencies. Concurrently, Apache Arrow \cite{arrow2024} defined a standardized in-memory columnar format, while Velox \cite{pedreira2022velox} provided a modular, high-performance C++ vectorized execution library deployed across modern data infrastructure. In distributed computing, Photon \cite{behm2022photon} demonstrated the efficacy of vectorized execution kernels for large-scale lakehouse analytical processing, while learned query optimization techniques \cite{vldb2023queryopt} have demonstrated potential in tuning analytical execution plans. More recently, Polars \cite{polars2024} introduced a multithreaded DataFrame engine written in Rust that incorporates lazy query evaluation and streaming execution.

In contrast to traditional in-memory Python libraries such as Pandas \cite{pandas2023}, which materialize complete, uncompressed matrices in heap memory, vectorized columnar engines support streaming execution pipelines. Hexanna IIVX leverages these vectorized foundations by embedding an out-of-core execution engine as an architectural adapter, bridging high-level cognitive agent reasoning with vectorized execution kernels.

\subsection{Out-of-Core Processing and Memory-Bounded Query Evaluation}
Out-of-core query processing remains a foundational pillar of database systems research, particularly under modern high-bandwidth solid-state architectures \cite{bandle2022nvme}. When relational operations (such as hash aggregations, sorted groupings, or multi-way joins) demand more buffer memory than physically allocated, out-of-core engines partition intermediate data into disk spill blocks, swapping partitions between volatile memory and secondary NVMe storage \cite{kohn2022adaptive}. Dreseler et al. \cite{dreseler2023benchmarking} quantified memory consumption trade-offs across modern database engines, highlighting that controlled disk spilling prevents operating system crashes at the cost of disk I/O latency. The TPC-H decision support benchmark \cite{tpch2021} provides the industry-standard benchmark suite to evaluate system behavior under complex multi-table joins, subqueries, and intensive aggregations across varying Scale Factors.

While modern analytical engines implement internal spilling mechanisms, autonomous agents historically have operated oblivious to memory ceilings. Agents that synthesize raw Pandas scripts frequently fail when handling multi-gigabyte workloads because Pandas lacks out-of-core spilling primitives. Hexanna IIVX explicitly operationalizes out-of-core query processing within an autonomous agent harness, enforcing strict buffer limits to ensure guaranteed execution stability.

\subsection{LLM Security, Execution Sandboxing, and Prompt Injections}
Deploying autonomous LLM agents with code execution capabilities introduces severe operational and security vulnerabilities. Toolformer \cite{schick2023toolformer} demonstrated that language models can autonomously invoke external tools to solve complex tasks. However, when agents are granted unrestricted shell or operating system access, they become susceptible to indirect prompt injection attacks \cite{greshake2023not}. Malicious payloads embedded within untrusted data files can hijack agent execution flows, leading to data exfiltration or host system compromise \cite{liu2023jailbreaking}.

Recent empirical evaluations by Li et al. \cite{li2024evaluating} underscore that software sandboxing boundaries in LLM agents are frequently bypassed when interpreters expose low-level system calls or arbitrary execution environments (e.g., Python \texttt{eval()} or Bash subshells). Even in multi-agent and simulated environments \cite{park2023generative}, unrestricted action spaces produce non-deterministic side effects. Hexanna IIVX eliminates these vulnerability vectors by replacing unrestricted interpreters with a deterministic, strictly-typed functional tool harness. By confining agent operations to schema inspection and analytical SQL queries executed inside an isolated embedded engine, Hexanna IIVX enforces strong defense-in-depth against prompt injection and unauthorized system modification.
